%% Experimental Setup Section — updated 2026-06-17 to reflect QUASAR v15 + SiliQun
\section{Experimental Setup}
\label{sec:experiments}

\subsection{QUASAR v15 Framework}

QUASAR v15 (Quantum Universal Adaptive State pReparation, version 15) is the
production implementation used for all experiments reported in this paper.
It extends QUASAR v14 with a six-layer modular architecture connected by a
typed Signal Bus, implementing a \emph{proactive-plus-reactive} training
philosophy: the outer DEHB loop establishes an optimal hyperparameter
configuration before training begins, while the reactive pipeline (SWDFT
$\to$ ERL $\to$ SDFT $\to$ TTT $\to$ DEHB inner) provides mid-flight
correction without restarting.

\paragraph{Base RL algorithm.}
We use Soft Actor-Critic (SAC) \cite{haarnoja2018soft} with automatic entropy
tuning and Feature-wise Linear Modulation (FiLM) \cite{perez2018film} for
goal conditioning. The actor and critic networks are multilayer perceptrons
with two hidden layers of 256 units each and ReLU activations, with FiLM
layers injecting the target state embedding $\phi_{\mathrm{tgt}} \in \mathbb{R}^{2^n}$
at each hidden layer. The replay buffer has capacity $10^6$ transitions and
we use a batch size of 256.

\paragraph{Continual multi-target learning.}
To enable a single agent to learn control policies for multiple entangled state
families without catastrophic forgetting, we augment SAC with Dark Experience
Replay++ (DER++) \cite{buzzega2020dark}. DER++ stores a buffer of (state, logit)
pairs from past training trajectories and adds a distillation loss term to the
SAC objective:
\begin{equation}
  \mathcal{L}_{\mathrm{DER++}} = \mathcal{L}_{\mathrm{SAC}} +
  \lambda \mathbb{E}_{(s, z) \sim \mathcal{B}_{\mathrm{DER}}}
  \bigl[\|Q(s, \cdot) - z\|^2\bigr],
\end{equation}
where $\mathcal{B}_{\mathrm{DER}}$ is the DER++ buffer and $\lambda = 0.5$.

\paragraph{Adaptive Convergence Control.}
The Adaptive Convergence Controller (ACC) replaces fixed training budgets and
heuristic plateau counters with a principled stopping criterion based on
saturating exponential model fitting (see Section~\ref{sec:acc}). In v15,
ACC is integrated with the Stochastic Layer Modulation (SLM) module: SLM
fires only when the Signal Bus receives a \textsc{Stagnation} signal from the
SWDFT detector, not on a fixed timer. This makes SLM a reactive corrector
rather than a periodic one, reducing wasted compute on already-converging runs.

\paragraph{Stagnation detection and reactive correction.}
The Sliding-Window Discrete Fourier Transform (SWDFT) detector monitors the
fidelity trajectory $\{F_t\}$ using a window of 2,000 steps. Stagnation is
declared when the dominant spectral power falls below a threshold $\gamma$,
indicating that the trajectory has become stationary. Upon stagnation, the
Emergency Reactive Layer (ERL) escalates through three correction stages:
\begin{enumerate}
  \item \textbf{ERL / SLM (Stage 1):} BRFD-calibrated probe episodes (200--500
    steps) with supervised actor update. Cost: $O(500)$ environment steps.
  \item \textbf{SDFT (Stage 2):} Spectral decomposition of the fidelity
    trajectory to identify the dominant plateau frequency, followed by a
    targeted perturbation of the critic's value estimates.
    Cost: $O(2{,}000)$ environment steps.
  \item \textbf{TTT (Stage 3):} Test-time training on the current observation
    distribution using the historical best policy checkpoint as a teacher.
    Cost: $O(5{,}000)$ environment steps.
\end{enumerate}
If all three stages fail to escape the plateau, the inner DEHB loop performs a
narrow re-search over $\alpha$ and $\mathrm{lr\_actor}$ (8 brackets, 10,000
steps each) to escape at the hyperparameter level.

\paragraph{Bayesian reward function discovery.}
The BRFD module learns the reward weight vector
$\mathbf{w} = (w_F, w_S, w_{\mathrm{str}}, w_{\log})$ jointly with the
policy using a Gaussian process surrogate over the reward parameter space.
BRFD calibrates the probe budget for each ERL invocation based on the
current training difficulty, as measured by the ACC model's predicted
$F_\infty$ and $\tau$ parameters.

\paragraph{Hyperparameter optimisation.}
All architecture-wide hyperparameters are learned by the outer DEHB loop
\cite{awad2021dehb} with 36 successive halving brackets. Each bracket runs a
short-budget trial (40,000 steps), with ACC providing early termination for
unreachable configurations. The best configuration from the HP search phase
is used for the full-budget run, with the budget set to $T^* \times 1.2$ as
predicted by ACC's saturating exponential model.

\subsection{Quantum Simulation Environment}

We simulate $n$-qubit systems using SiliQun v2 \cite{TODO_siliqun}, a
modular silicon spin qubit simulator that models exchange-coupled spin-1/2
particles with realistic device parameters. SiliQun supports three simulation
backends: Matrix Product State (MPS) for pure-state simulation (used for
$n < 8$), and GPU-accelerated exact state-vector simulation (used for
$n \geq 8$). The Hamiltonian is:
\begin{equation}
  H_0 = \sum_{i=1}^{n} \frac{\omega_i}{2} \sigma_z^{(i)}
        + \sum_{\langle i,j\rangle} J_{ij} \bigl(\sigma_x^{(i)}\sigma_x^{(j)}
        + \sigma_y^{(i)}\sigma_y^{(j)}\bigr),
\end{equation}
where $\omega_i / 2\pi \in [4.5, 5.5]\,\mathrm{GHz}$ are qubit frequencies and
$J_{ij} / 2\pi \in [1, 10]\,\mathrm{MHz}$ are exchange coupling strengths,
drawn from the donor device profile \cite{TODO_siliqun}. The action space
consists of parameterised single-qubit rotations and two-qubit gates with
continuous angle parameters in $[-1, 1]$.

\subsection{Target State Families}

We evaluate QUASAR v15 on four canonical entangled state families across
qubit counts $n \in \{2, 3, \ldots, 12\}$:
\begin{enumerate}
  \item \textbf{GHZ states:} $\ket{\mathrm{GHZ}_n} = (\ket{0}^{\otimes n} + \ket{1}^{\otimes n})/\sqrt{2}$.
  \item \textbf{W states:} $\ket{W_n} = n^{-1/2}\sum_{i=1}^{n}\ket{0\cdots 1_i \cdots 0}$.
  \item \textbf{Linear cluster states:} $\ket{\mathrm{CL}_n}$, prepared by
    applying $H^{\otimes n}$ followed by $CZ$ on all neighbouring pairs.
  \item \textbf{Dicke states:} $\ket{D_n^k}$ with $k = 3$ excitations,
    equal superposition of all $\binom{n}{k}$ computational basis states
    with Hamming weight $k$.
\end{enumerate}

\subsection{Baseline Methods}

We compare ACC against three baselines:
\begin{enumerate}
  \item \textbf{Fixed budget (FB)}: the standard approach used in prior work
    \cite{yao2021reinforcement,moro2021quantum}, with a fixed budget of
    $500{,}000 \times 1.5^{n-2}$ steps per qubit level.
  \item \textbf{Plateau counter (PC)}: the heuristic used in QUASAR v13,
    which stops training when the best fidelity has not improved for 20,000
    consecutive steps.
  \item \textbf{ASHA}: the Asynchronous Successive Halving Algorithm
    \cite{li2020system} applied to the training budget allocation problem,
    using the same 36-trial budget as DEHB+ACC.
\end{enumerate}

\subsection{Evaluation Protocol}

All experiments are run with three random seeds (42, 123, 456) per
(target family, qubit count) combination, yielding $4 \times 11 \times 3 = 132$
independent training runs. We report the mean and standard deviation of the
best fidelity achieved, the total number of training steps consumed, and the
wall-clock time on the Aziz Supercomputer A100 nodes. Statistical significance
is assessed using a two-sided Wilcoxon signed-rank test with $\alpha = 0.05$.

\subsection{Computational Resources}

All experiments were performed on the Aziz Supercomputer at King Abdulaziz
University, using NVIDIA A100-PCIE-40GB GPUs. Each seed is assigned a
dedicated A100 node (jobs 181560, 181561, 181562 for seeds 42, 123, 456
respectively), running in parallel under PBS Pro with a 72-hour walltime limit.
The total compute budget for the main experiments is approximately 216 GPU-hours
($3 \times 72\,\mathrm{h}$).
